%% Pattern Recognition (Elsevier) 投稿模板
%% 文档类: elsarticle (TeX Live 自带, 或 CTAN/Overleaf)
%% [review] = 双倍行距+行号(投稿审稿用); [3p] = 单栏; 最终录用改 [final,5p]
%% 投稿前必读: https://www.sciencedirect.com/journal/pattern-recognition/publish/guide-for-authors
\documentclass[review,3p]{elsarticle}

\usepackage{lineno}
\modulolinenumbers[5]

\usepackage{graphicx}
\usepackage{amsmath,amssymb,amsfonts}
\usepackage{booktabs}
\usepackage{multirow}
\usepackage{algorithm}
\usepackage{algpseudocode}
\usepackage{hyperref}
\usepackage{xcolor}

\journal{Pattern Recognition}

\begin{document}

\begin{frontmatter}
\title{Iris-Based Kinship Recognition in Pigeons via\\
Pedigree-Graph Supervised Representation Learning}

\author[a]{First Author}
\author[a]{Second Author\corref{cor1}}
\cortext[cor1]{Corresponding author}

\affiliation[a]{organization={School/Department, University},
                addressline={Address line},
                city={City},
                postcode={000000},
                country={China}}

\begin{abstract}
We study a novel problem: \emph{kinship recognition from iris images in pigeons}---determining the genealogical relatedness of two pigeons solely from their iris texture. We release a dataset of 25{,}690 normalized iris images linked to 126{,}105 pedigree records extracted from raw bloodline-text data. We make two core contributions. (1) We show that iris feature distance correlates strongly with true pedigree kinship (Spearman $-0.716$, AUC $0.93$ on unseen individuals), demonstrating iris carries bloodline information. (2) We define a graded kinship metric from a pedigree graph built by parsing structured father/mother/grandparent fields, combining a literal ancestor-path code (recursive parent-code concatenation + LCS) with a founder-contribution vector whose dot product yields the kinship coefficient. We further propose \emph{pedigree-supervised graded SupCon}, training the iris encoder with the pedigree kinship as graded relevance, which improves graded nDCG@20 by 19\% over the IDF-heuristic proxy while restoring fine-grained tier monotonicity. Code and data are released.
\end{abstract}

\begin{keyword}
Iris recognition \sep Kinship verification \sep Representation learning \sep
Metric learning \sep Pedigree graph \sep Supervised contrastive learning
\end{keyword}
\end{frontmatter}

\linenumbers

%% ============ 正文(内容见 paper/paper_draft.md, 逐节填入) ============

\section{Introduction}
\label{sec:intro}
% 见 paper_draft.md 第1节. 要点: 信鸽育种血统验证需求; 虹膜能否判血缘的假设;
% 与人虹膜识别(身份非血缘)、人脸血缘验证(不同模态)的区别; 三贡献.

\section{Related Work}
\label{sec:related}
% 2.1 Iris recognition (Daugman, 人虹膜身份; 无鸽虹膜/血缘前作)
% 2.2 Kinship verification (人脸血缘验证, 多为二值; 本文扩展到 graded multi-label)
% 2.3 Metric learning (SupCon, ArcFace, Proxy-Anchor; graded-relevance weighted SupCon)
% 2.4 Pedigree/genetic kinship (numerator relationship matrix, Henderson; kinship coefficient $\phi$)

\section{Method}
\label{sec:method}
\subsection{Pedigree Graph Construction}
% 解析 details.txt (父亲/母亲/祖父/祖母/外祖父/外祖母:足环号); 53,158 节点 / 19,642 边 /
% 42,609 founder; 环号合法性过滤; founder=无入度; 5 记录环用路径内 visited 处理.

\subsection{Kinship Encoding}
% 形式A(字面): founder 唯一码 + 后代=父母码递归拼接(无界到 founder) + LCS 相似度.
% 形式B(向量): $v_i[a]=\sum_{\text{paths }i\to a}(1/2)^{\text{depth}}$, $k(i,j)=v_i\cdot v_j \approx 2\phi_{ij}$.
% Hybrid: 无结构系谱者用 IDF 加权 blood_id 集合兜底.

\subsection{Pedigree-Supervised Graded SupCon}
% Relevance matrix $R_{ij}=k(i,j)$ (hybrid). Loss:
% \[ L = -\log\frac{\sum_j \exp(s_{ij}/\tau)\,R_{ij}}{\sum_j \exp(s_{ij}/\tau)\,(R_{ij}^{+} + (1-R_{ij}^{+}))} \]
% $R_{ij}^{+}=R_{ij}\,\mathbf{1}[R_{ij}\ge\text{cutoff}]$ (弱亲缘当负样本保 tier);
% fallback 降权(fallback\_scale=0.3)使 pedigree k 主导.

\section{Dataset}
\label{sec:dataset}
% 25,690 归一化虹膜图(64x512, Daugman+mask); 31,896 原始图;
% pigeon.csv(PG_ID/BLOOD) + relations.csv(250,207 image-ancestor 对) + details.txt(系谱);
% clean split 17,983/2,526; 局限: 结构系谱覆盖14.8%, 64%孤立, 记录推导非基因.

\section{Experiments}
\label{sec:exp}
\subsection{Iris distance vs pedigree kinship (Contribution 1)}
% Table: Spearman -0.716, AUC full/half/cousin 0.954/0.946/0.906, 单调.
% Fig: 散点(iris_dist vs k) + 分层箱线.

\subsection{Kinship encoding validation (Contribution 2)}
% Table: 全同胞0.715>半同胞0.400>表亲0.104>无关0.000, 单调, 与理论吻合.

\subsection{Pedigree-supervised training (Contribution 3)}
% Table (公平 150ep): idf -0.711/0.467/0.921; hybrid(v2) -0.716/0.558/0.931; ped_only -0.601/0.280/0.882(FAIL).

\subsection{Ablations}
% D1 kinship source (上表); D2 depth; D3 graded vs binary; D4 cutoff sweep.

\section{Conclusion}
\label{sec:conclusion}
% 信鸽虹膜血缘识别新问题; 数据集+编码+方法三贡献; 局限与 broader impact (育种防伪/误用).

\section*{Declaration of competing interest}
The authors declare no competing financial or personal interests.

\section*{Acknowledgments}
% 基金/致谢.

%% 参考文献: Pattern Recognition 用 numbered 风格 (elsarticle-num); 也接受 author-year.
%% 投稿前在 Guide for Authors 确认; refs 见 refs.bib
\bibliographystyle{elsarticle-num}
\bibliography{refs}

\end{document}
